\documentclass[compress,12pt]{beamer}

\usetheme{Arguelles}

\usepackage{graphicx}
\usepackage{caption}
\usepackage[spanish,es-noshorthands]{babel}
\usepackage[pages=some]{background}

\usepackage{tikz}
\usepackage{tikz-cd}
\usetikzlibrary{
  babel,
  arrows.meta,
  calc,
  positioning,
  overlay-beamer-styles,
  shadows.blur
}
\usetikzlibrary{decorations.pathmorphing}

\usepackage{amsmath,amssymb,latexsym,amscd} 
\usepackage[all,cmtip]{xy}
\usepackage{fancyhdr}
\usepackage{mathalfa}
\usepackage{mathrsfs}
\usepackage{hyperref}
\usepackage{ragged2e}
\usepackage{wasysym}
\usepackage[authoryear]{natbib}

% Colores para la diapositiva Top/Frm
\definecolor{TopBlue}{RGB}{14,73,160}
\definecolor{FrmGreen}{RGB}{15,120,45}
\definecolor{TopLight}{RGB}{235,242,252}
\definecolor{FrmLight}{RGB}{235,248,238}

\pagenumbering{arabic}

\DeclareMathOperator{\op}{op}
\DeclareMathOperator{\pt}{pt}
\DeclareMathOperator{\tp}{tp}
\DeclareMathOperator{\Sp}{sp}
\DeclareMathOperator{\SP}{Sp}
\DeclareMathOperator{\Sob}{Sob}
\DeclareMathOperator{\spec}{spec}
\DeclareMathOperator{\Fit}{Fit}
\DeclareMathOperator{\Pth}{P}
\DeclareMathOperator{\Vrm}{Vrm}
\DeclareMathOperator{\Frm}{Frm}
\DeclareMathOperator{\Top}{Top}
\DeclareMathOperator{\Ord}{Ord}
\DeclareMathOperator{\Obj}{Obj}
\DeclareMathOperator{\Hom}{Hom}

\title{De espacios a marcos:}
\subtitle{la teoría de categorías como puente hacia la topología sin puntos}
\event{CIMA 2026, BUAP}
\date{\today}
\author{Juan Carlos Monter Cortés}
\institute{Universidad de Guadalajara}
\email{juan.monter2902@alumnos.udg.mx}

%\homepage{www.mywebsite.com}
\github{JCmonter}

\begin{document}


%\frame[plain]{\titlepage}

\begin{frame}[plain]
\titlepage

\vfill
\begin{center}
    \textcolor{white}{\scriptsize Investigación apoyada por SECIHTI-PROYECTO CBF2023-2024-2630}
\end{center}
\end{frame}

%\begin{frame}{Contenido}
%\tableofcontents
%\end{frame}

\begin{frame}{¿Cuánta información de un espacio topológico está contenida en sus abiertos?}
\small{
\begin{columns}[c]

% Columna izquierda
\begin{column}{0.38\textwidth}
\centering

\[
X=\{x,y,z\}
\]

\vspace{0.3cm}

\begin{tikzpicture}[scale=1.1]
    \fill (0,0) circle (2pt) node[below] {$x$};
    \fill (1.2,0.8) circle (2pt) node[right] {$y$};
    \fill (2.2,-0.1) circle (2pt) node[below] {$z$};
\end{tikzpicture}
\end{column}

% Columna derecha
\begin{column}{0.58\textwidth}

\[
\mathcal O(X)
=
\left\{
\varnothing,\,
\{z\},\,
\{y,z\},\,
X
\right\}.
\]

\centering

\begin{tikzpicture}[
    node distance=0.65cm,
    every node/.style={inner sep=1pt}
]
\node (X) {$X$};
\node[below=of X] (yz) {$\{y,z\}$};
\node[below=of yz] (z) {$\{z\}$};
\node[below=of z] (empty) {$\varnothing$};

\draw (X) -- (yz);
\draw (yz) -- (z);
\draw (z) -- (empty);
\end{tikzpicture}

\end{column}

\end{columns}
}

\[
X \longmapsto \mathcal O(X)
\]

\begin{block}{Pregunta}
¿Qué información topológica permanece si olvidamos los puntos y
 conservamos únicamente sus abiertos?
\end{block}

\end{frame}

\begin{frame}{La estructura algebraica de los abiertos}
\small{
Sea $X$ un espacio topológico. La colección $\mathcal O(X)$ de sus
abiertos posee naturalmente las siguientes operaciones:

\begin{columns}[T]

    \begin{column}{0.47\textwidth}
        \begin{block}{Supremos arbitrarios}
            Si $\{U_i\}_{i\in I}\subseteq\mathcal O(X)$, entonces
            \[
                \bigvee_{i\in I}U_i
                =
                \bigcup_{i\in I}U_i.
            \]
        \end{block}
    \end{column}

    \begin{column}{0.47\textwidth}
        \begin{block}{Ínfimos finitos}
            Si $U,V\in\mathcal O(X)$, entonces
            \[
                U\wedge V
                =
                U\cap V.
            \]
        \end{block}
    \end{column}

\end{columns}

Estas operaciones satisfacen la ley distributiva

\[
    U\wedge\bigvee_{i\in I}V_i
    =
    \bigvee_{i\in I}(U\wedge V_i).
\]

    \centering
    ¿Podemos tomar esta estructura algebraica como nuestro punto de partida?

}
\end{frame}

\begin{frame}{Marcos}
\small{
\begin{block}{Definición}
Un \textbf{marco} es un retículo completo $A$ que satisface la
ley distributiva
\[
    a\wedge\bigvee S
    =
    \bigvee_{s\in S}(a\wedge s),
    \qquad a\in A,\quad S\subseteq A.
\]
\end{block}

\pause

\medskip

El ejemplo fundamental proviene de la topología:
\[
\boxed{
    X\text{ espacio topológico}
    \qquad\Longrightarrow\qquad
    \mathcal O(X)\text{ es un marco}.
}
\]

\pause

\medskip

En $\mathcal O(X)$,
\[
    \bigvee_{i\in I}U_i=\bigcup_{i\in I}U_i,
    \qquad
    U\wedge V=U\cap V.
\]

\begin{center}
\textbf{La estructura de marco puede estudiarse sin hacer referencia
a los puntos de $X$.}
\end{center}
}
\end{frame}

\begin{frame}{¿Qué ocurre con las funciones continuas?}
\small{
Sea
\[
    f:X\longrightarrow Y
\]
una función continua.

\pause

Entonces la imagen inversa induce una aplicación
\[
    f^{-1}:\mathcal O(Y)\longrightarrow\mathcal O(X),
    \qquad
    V\longmapsto f^{-1}(V).
\]

\pause


Además, $f^{-1}$ preserva la estructura de marco:
\[
    f^{-1}\left(\bigcup_{i\in I}V_i\right)
    =
    \bigcup_{i\in I}f^{-1}(V_i),
\quad\text{y}\quad
    f^{-1}(U\cap V)
    =
    f^{-1}(U)\cap f^{-1}(V).
\]

\centering
Una función continua entre espacios induce un morfismo
entre sus marcos de abiertos, \textbf{pero en dirección contraria}.
}
\end{frame}

\begin{frame}{De espacios a marcos}
\small{
Hemos visto que
\[
    X \longmapsto \mathcal O(X)
\]
asocia a cada espacio topológico su marco de abiertos.

\pause

Además, si
\[
    f:X\longrightarrow Y
\]
es continua, entonces
\[
    f^{-1}:\mathcal O(Y)\longrightarrow\mathcal O(X)
\]
es un morfismo de marcos.

\pause

\medskip

Por lo tanto, obtenemos un funtor contravariante
\[
    \mathcal O:\mathbf{Top}\longrightarrow\mathbf{Frm}^{op}=\mathbf{Loc}.
\]
}
\end{frame}

\begin{frame}{¿Podemos recuperar los puntos?}
\small{
\begin{block}{Pregunta}
Dado un marco $A$, ¿podemos asociarle un espacio topológico
de manera natural?
\end{block}

\pause

\medskip

La respuesta consiste en considerar sus \textbf{puntos}.

\[
    A
    \quad\longmapsto\quad
    \operatorname{pt}(A).
\]

\pause

\vfill

De esta manera aparecen dos construcciones:

\[
\boxed{
    X\longmapsto\mathcal O(X)
    \qquad\text{y}\qquad
    A\longmapsto\operatorname{pt}(A).
}
\]
}
\end{frame}

\begin{frame}{Los puntos de un marco}
Sea $A$ un marco. Un elemento $p\in A$, con $p\neq 1$, es
\textbf{$\wedge$-irreducible} si

\[
    a\wedge b\leq p
    \quad\Longrightarrow\quad
    a\leq p \;\text{ o }\; b\leq p.
\]

\pause

\medskip

Denotamos por

\[
    \operatorname{pt}(A)
    =
    \{p\in A : p \text{ es } \wedge\text{-irreducible}\}
\]

al conjunto de puntos de $A$, equipado con la topología generada por los conjuntos abiertos
\[
    \mathcal{U}(a)=\{p\in \pt A\mid a\nleq p\}.
\]

\end{frame}

\begin{frame}[fragile]{La adjunción entre espacios y marcos}
\small{
Hemos construido dos asignaciones naturales:

\[
    X \longmapsto \mathcal O(X),
    \qquad
    A \longmapsto \operatorname{pt}(A).
\]

\pause

\medskip

Estas construcciones se relacionan mediante una adjunción:

\[
\begin{tikzcd}[column sep=huge]
    \mathbf{Top}
    \arrow[r, shift left=1.2ex, "\mathcal O"]
&
    \mathbf{Frm}^{op}
    \arrow[l, shift left=1.2ex, "\operatorname{pt}"]
\end{tikzcd}
\]

\[
    \mathcal O \dashv \operatorname{pt}.
\]

\pause

\begin{block}{}
\centering
La teoría de categorías proporciona el lenguaje natural
para comparar espacios topológicos y marcos.
\end{block}
}
\end{frame}

\begin{frame}{¿Cuándo recuperamos el objeto original?}
\small{
La adjunción nos permite recorrer ambos caminos:

\medskip

\begin{columns}[T]

% Columna izquierda
\begin{column}{0.47\textwidth}

\begin{block}{Desde un espacio}
\[
X
\longmapsto
\mathcal O(X)
\longmapsto
\operatorname{pt}(\mathcal O(X)).
\]

\medskip

Existe una aplicación natural
\[
X\longrightarrow
\operatorname{pt}(\mathcal O(X)).
\]

\medskip

Cuando ésta es un homeomorfismo,
decimos que $X$ es \textbf{sobrio}.
\end{block}

\end{column}

% Columna derecha
\begin{column}{0.47\textwidth}

\begin{block}{Desde un marco}
\[
A
\longmapsto
\operatorname{pt}(A)
\longmapsto
\mathcal O(\operatorname{pt}(A)).
\]

\medskip

Existe un morfismo natural
\[
A\longrightarrow
\mathcal O(\operatorname{pt}(A)).
\]

\medskip

Cuando éste es un isomorfismo,
decimos que $A$ es \textbf{espacial}.
\end{block}

\end{column}

\end{columns}

\pause

\vfill

\begin{center}
\boxed{
\mathbf{Sob}
\;\simeq\;
\mathbf{Sp}^{op}
}
\end{center}
}
\end{frame}

\begin{frame}{¿Y si el marco no es espacial?}
\small{
Existen marcos $A$ para los cuales

\[
A \not\cong \mathcal O(X)
\]

para ningún espacio topológico $X$.

\pause

\bigskip

\begin{block}{Topología sin puntos}
La idea es estudiar la estructura de marco como un objeto
topológico por derecho propio, sin exigir que sus elementos
sean los abiertos de un espacio de puntos.
\end{block}

\pause

\vfill

\begin{center}
¿Qué conceptos y construcciones topológicas sobreviven
sin utilizar puntos?
\end{center}
}
\end{frame}

\begin{frame}{Axiomas de separación}
\small
En topología clásica, muchas propiedades se formulan
directamente en términos de puntos.

\medskip

Por ejemplo, un espacio $X$ es \textbf{$T_1$} si

\[
    x\neq y
    \quad\Longrightarrow\quad
    \text{existe } U\in\mathcal O(X)
    \text{ tal que }
    x\in U,\; y\notin U.
\]

\pause

\medskip

Equivalentemente,

\[
    X \text{ es } T_1
    \quad\Longleftrightarrow\quad
    \{x\} \text{ es cerrado para todo }x\in X.
\]

\pause

\vfill

\begin{block}{El problema}
\centering
Si queremos trabajar únicamente con un marco $A$,
¿cómo formulamos una propiedad que originalmente
habla de \textbf{puntos}?
\end{block}

\end{frame}

\begin{frame}{Reformular antes de generalizar}
\small

Para trasladar una propiedad topológica al contexto de marcos
buscamos una formulación que dependa únicamente de los abiertos.

\medskip

\begin{center}
\begin{tikzpicture}[
    node distance=0.75cm,
    box/.style={
        draw,
        rounded corners,
        align=center,
        text width=2.5cm,
        minimum height=0.9cm,
        inner sep=3pt
    }
]

\node[box] (top) {Propiedad\\topológica};

\node[box, right=of top] (open)
    {Formulación \\en términos\\ de abiertos};

\node[box, right=of open] (frame)
    {Propiedad\\del marco};

\draw[->] (top) -- (open);
\draw[->] (open) -- (frame);

\end{tikzpicture}
\end{center}

\pause

\medskip

Así, una propiedad inicialmente formulada sobre un espacio $X$
\[
    \text{puede adquirir sentido directamente sobre } \mathcal O(X),
\]

y posteriormente sobre un marco arbitrario $A$.

\medskip

\begin{block}{}
\centering
La traducción no consiste en eliminar los puntos de la notación,
sino en encontrar la estructura algebraica que expresa
la misma información topológica.
\end{block}

\end{frame}

\begin{frame}{Un ejemplo: regularidad}
\small

Recordemos que un espacio topológico $X$ es \textbf{regular} si,
dados un punto $x\in X$ y un cerrado $C$ con $x\notin C$,
existen abiertos disjuntos que los separan.

\pause

\medskip

Esta propiedad puede expresarse únicamente en términos de abiertos.

\begin{block}{La relación \textit{well inside}}
Para $U,V\in\mathcal O(X)$, escribimos $U\prec V$ si existe $W\in\mathcal O(X)$ tal que
\[
    U\cap W=\varnothing
    \qquad\text{y}\qquad
    V\cup W=X.
\]
\end{block}

\pause

Entonces la regularidad puede expresarse como

\[
    U
    =
    \bigcup\{V\in\mathcal O(X):V\prec U\},
    \qquad U\in\mathcal O(X).
\]

\end{frame}

\begin{frame}{Olvidemos el espacio}
\small
Para $a,b\in A$, definimos

\[
    a\prec b
    \quad\Longleftrightarrow\quad
    \exists\,c\in A
    \text{ tal que }
    a\wedge c=0
    \quad\text{y}\quad
    b\vee c=1.
\]

\pause

\medskip

\begin{block}{Marco regular}
Un marco $A$ es \textbf{regular} si, para todo $a\in A$,
\[
    a=\bigvee\{b\in A:b\prec a\}.
\]
\end{block}

\pause

\vfill

\begin{center}
\boxed{
\text{La definición ya no necesita un espacio de puntos.}
}
\end{center}

\end{frame}

\begin{frame}{De propiedades a construcciones}

Hasta ahora hemos visto cómo una \textbf{propiedad topológica}
puede reformularse en el lenguaje de marcos:

\[
\boxed{
\text{Espacios}
\longrightarrow
\text{Abiertos}
\longrightarrow
\text{Marcos}
}
\]

\pause

\medskip

Pero podemos hacer la misma pregunta para
\textbf{construcciones topológicas}.

\pause

\bigskip

\begin{center}
\Large
¿Qué ocurre, por ejemplo, con la modificación \textit{patch}?
\end{center}

\end{frame}

\begin{frame}{La modificación \textit{patch}}
\small

Sea $X$ un espacio topológico.

\medskip

La idea de una modificación \textit{patch} es considerar
una topología más fina sobre el mismo conjunto de puntos:

\[
\boxed{
(X,\tau)
\qquad\longmapsto\qquad
(X,\tau_{\mathrm{patch}})
}
\]

\pause

\medskip

Se conservan los puntos de $X$, pero se incorporan nuevos abiertos
a partir de información que ya estaba presente en la topología original.

\pause

\bigskip

Por tanto, $\tau\subseteq\tau_{\mathrm{patch}}$.


\begin{block}{Idea}
\centering
El espacio no cambia como conjunto; lo que cambia es
la manera en que observamos topológicamente sus puntos.
\end{block}

\end{frame}

\begin{frame}{¿Cómo hacemos \textit{patch} sin puntos?}
\small

En espacios, una modificación \textit{patch} cambia la topología:

\[
(X,\tau)
\longmapsto
(X,\tau_{\mathrm{patch}}).
\]

\pause

\medskip

Pero si nuestro objeto de partida es un marco $A$,
ya no tenemos que comenzar con un conjunto de puntos.

\bigskip

\begin{center}
\begin{tikzpicture}[
    node distance=1.3cm,
    box/.style={
        draw,
        rounded corners,
        align=center,
        minimum width=2.7cm,
        minimum height=1cm
    }
]

\node[box] (space) {Espacio\\$X$};
\node[box, right=of space] (opens) {Marco de abiertos\\$\mathcal O(X)$};
\node[box, right=of opens] (patch) {Construcción\\\textit{patch}};

\draw[->] (space) -- (opens);
\draw[->] (opens) -- (patch);

\end{tikzpicture}
\end{center}

\pause

\bigskip

\begin{block}{Pregunta}
\centering
¿Podemos describir la modificación \textit{patch}
utilizando únicamente la estructura del marco?
\end{block}

\end{frame}

\begin{frame}{El \textit{patch} en el mundo point-free}
\small

La construcción \textit{patch} también admite una formulación
intrínseca en el lenguaje de marcos.

\pause

\medskip

En lugar de modificar directamente los puntos de un espacio,
trabajamos con la estructura interna del marco:

\[
\boxed{
A
\qquad\longmapsto\qquad
\operatorname{Patch}(A)
}
\]

\pause

\medskip

Para ello aparecen herramientas propias de la teoría de marcos,
como los \textbf{núcleos}, que permiten describir sublocales
y construir nuevas estructuras a partir de $A$.

\pause

\vfill

\begin{block}{}
\centering
Una construcción que nació en términos de espacios
puede continuar teniendo sentido aun cuando
el marco no sea espacial.
\end{block}

\end{frame}

\begin{frame}{¿Qué ganamos con esta perspectiva?}

\begin{columns}[T]

% Primera columna
\begin{column}{0.30\textwidth}
\centering

\Large
\textbf{Traducir}

\normalsize
\medskip

\[
\mathbf{Top}
\longleftrightarrow
\mathbf{Frm}
\]

\medskip

Interpretar propiedades y construcciones topológicas
en términos algebraicos.

\end{column}

\pause

% Segunda columna
\begin{column}{0.30\textwidth}
\centering

\Large
\textbf{Generalizar}

\normalsize
\medskip

\[
\mathcal O(X)
\rightsquigarrow
A
\]

\medskip

Extender conceptos desde marcos espaciales
hacia marcos arbitrarios.

\end{column}

\pause

% Tercera columna
\begin{column}{0.30\textwidth}
\centering

\Large
\textbf{Descubrir}

\normalsize
\medskip

\[
\text{¿?}
\]

\medskip

Encontrar nuevos fenómenos y preguntas
que no dependen de la existencia de puntos.

\end{column}

\end{columns}

\vfill

\begin{block}{}
\centering
La traducción categórica no es únicamente un cambio de lenguaje:
permite ampliar el dominio de la topología.
\end{block}

\end{frame}

\begin{frame}{¿Qué preguntas aparecen?}
\small

\begin{block}{Separación}
¿Qué axiomas de separación admiten una formulación
puramente algebraica en marcos?
\end{block}

\pause

\begin{block}{Conservatividad}
Si una propiedad se define en el contexto point-free,
¿recuperamos la propiedad topológica original cuando
\[
A=\mathcal O(X)?
\]
\end{block}

\pause

\begin{block}{Más allá de los espacios}
¿Qué nuevos fenómenos aparecen cuando trabajamos con
marcos no espaciales?
\end{block}

\end{frame}

\begin{frame}{Para llevar a casa}

\begin{center}

\begin{tikzpicture}[
    node distance=1.6cm,
    box/.style={
        draw,
        rounded corners,
        align=center,
        minimum width=3cm,
        minimum height=1.2cm
    }
]

\node[box] (top) {
    Espacios topológicos\\
    $\mathbf{Top}$
};

\node[box, right=of top] (loc) {
    Marcos / Locales\\
    $\mathbf{Frm}^{op}=\mathbf{Loc}$
};

\draw[->, bend left=15, thick]
(top) to node[above] {$\mathcal O$} (loc);

\draw[->, bend left=15, thick]
(loc) to node[below] {$\operatorname{pt}$} (top);

\end{tikzpicture}

\medskip

\begin{minipage}{0.78\textwidth}
    \begin{block}{}
        \centering
        La teoría de categorías proporciona el puente
        entre ambos mundos.
    \end{block}
\end{minipage}

\pause

\bigskip

Los abiertos no solamente describen un espacio:

\medskip

{\Large
\textbf{pueden convertirse en el objeto de estudio.}
}

\end{center}

\end{frame}

\End

\begin{frame}{Referencias}
\small

\begin{itemize}

\item
J. Picado and A. Pultr,
\textit{Frames and Locales: Topology without Points},
Frontiers in Mathematics, Birkh\"auser, 2012.

\medskip

\item
J. Picado and A. Pultr,
\textit{Separation in Point-Free Topology},
Birkh\"auser, 2021.

\medskip

\item
P. T. Johnstone,
\textit{Stone Spaces},
Cambridge Studies in Advanced Mathematics,
Cambridge University Press, 1982.

\medskip

\item
R. A. Sexton,
\textit{A Point-Free and Point-Sensitive Analysis of the Patch Assembly},
Ph.D. Thesis, University of Manchester, 2003.

\end{itemize}

\end{frame}

%\section*{\textsc{Referencias}}
%\begin{frame}[allowframebreaks]
%\bibliographystyle{plainnat}
%\bibliographystyle{plain}
%\bibliographystyle{amsalpha}
%\bibliography{research2}
%\end{frame}

\begin{frame}[plain,standout]
      \centering
      \Huge{\bfseries \smiley¡Muchas gracias!\smiley}
      %In combination with \textit{plain}, \\
      %it makes a nice thank-you slide!
      %\scalebox{4}{\faGithub} \par\bigskip
      %\url{https://github.com/JCmonter/Apuntes/tree/main/Presentaciones} \\
      %\url{https://ctan.org/pkg/beamertheme-arguelles}
\end{frame}
\end{document}